% --- CORRECCIÓN IMPORTANTE EN LA PRIMERA LÍNEA ---
% Agregamos "xcolor=table" para que funcione \rowcolor sin errores
\documentclass[aspectratio=169, 10pt, xcolor=table]{beamer}

% --- Configuración del Tema (estilo profesional como el ejemplo) ---
\usetheme{Madrid}
\usecolortheme{dolphin}

% --- Paquetes necesarios ---
\usepackage[utf8]{inputenc}
\usepackage[spanish]{babel}
\usepackage{amsmath, amssymb}
\usepackage{booktabs}
\usepackage{graphicx}
\usepackage{listings}
\usepackage{tikz}
\usetikzlibrary{arrows.meta, positioning, shapes.geometric}

% --- Ruta de Imágenes (crea la carpeta "imagenes" y coloca ahí tus archivos) ---
\graphicspath{{./imagenes/}}

% --- Estilo para código Python ---
\definecolor{codegreen}{rgb}{0,0.6,0}
\definecolor{codegray}{rgb}{0.5,0.5,0.5}
\definecolor{codepurple}{rgb}{0.58,0,0.82}
\definecolor{backcolour}{rgb}{0.95,0.95,0.92}
\lstset{
    language=Python,
    backgroundcolor=\color{backcolour},
    commentstyle=\color{codegreen},
    keywordstyle=\color{blue},
    numberstyle=\tiny\color{codegray},
    stringstyle=\color{codepurple},
    basicstyle=\ttfamily\scriptsize,
    breaklines=true,
    frame=single,
    showstringspaces=false
}

% --- Pie de página personalizado ---
\makeatletter

\setbeamertemplate{footline}{
  \leavevmode%
  \hbox{%
  \begin{beamercolorbox}[wd=.333333\paperwidth,ht=2.25ex,dp=1ex,center]{author in head/foot}%
    \usebeamerfont{author in head/foot}\insertshortauthor
  \end{beamercolorbox}%
  \begin{beamercolorbox}[wd=.333333\paperwidth,ht=2.25ex,dp=1ex,center]{title in head/foot}%
    \usebeamerfont{title in head/foot}Sesión 1
  \end{beamercolorbox}%
  \begin{beamercolorbox}[wd=.333333\paperwidth,ht=2.25ex,dp=1ex,right]{date in head/foot}%
    \usebeamerfont{date in head/foot}Machine Learning \hfill \insertframenumber/\inserttotalframenumber \hspace*{0.3cm}
  \end{beamercolorbox}}%
  \vskip0pt%
}



\makeatother

% --- Datos de la portada ---
% --- Datos de la portada ---
\title[SESIÓN 1]{\textbf{Sesión 1}}  
\subtitle{Introducción al Machine Learning}
\author[Prof. C. Sánchez]{Carlos César Sánchez Coronel}
\institute{\textbf{MACHINE LEARNING}} 
\date{}


% Logo opcional (descomentar si tienes la imagen)
% \titlegraphic{\includegraphics[height=1.5cm]{logo.png}}



% (Opcional) Incluir un logo institucional, por ejemplo:
% \titlegraphic{\includegraphics[height=1.5cm]{logo.png}}

% --- Eliminar la barra de navegación (opcional) ---
\setbeamertemplate{navigation symbols}{}

% --- Índice al inicio de cada sección ---
\AtBeginSection[]{
  \begin{frame}
    \frametitle{Contenido}
    \tableofcontents[currentsection]
  \end{frame}
}




% ============================================================================
% INICIO DEL DOCUMENTO
% ============================================================================
\begin{document}

% --- Portada ---
\begin{frame}
    \titlepage
\end{frame}

% --- Índice general ---
\begin{frame}{Contenido}
    \tableofcontents
\end{frame}

% ============================================================================
% SECCIÓN 1: INTRODUCCIÓN Y TIPOS DE APRENDIZAJE
% ============================================================================
\section{Introducción y Tipos de Aprendizaje}

\begin{frame}{Inteligencia Artificial vs Machine Learning vs Deep Learning}
    % Basado en Pregunta 1 del solucionario
    \begin{block}{Jerarquía de conceptos}
        \begin{itemize}
            \item \textbf{Inteligencia Artificial (IA):} Campo amplio que busca crear sistemas con capacidades humanas (razonamiento, percepción, aprendizaje).
            \item \textbf{Machine Learning (ML):} Subcampo de la IA; algoritmos que aprenden de datos sin programación explícita.
            \item \textbf{Deep Learning (DL):} Subcampo del ML que usa redes neuronales profundas para modelar patrones complejos.
        \end{itemize}
        \vspace{0.3cm}
        \centering
        IA $\supset$ ML $\supset$ DL
    \end{block}
    % IMAGEN: Diagrama de Venn o jerárquico
\end{frame}

\begin{frame}{Tipos de Aprendizaje}
    % Basado en Pregunta 2 del solucionario
    \begin{columns}[t]
        \column{0.48\textwidth}
        \begin{exampleblock}{Supervisado}
            Datos etiquetados. \\
            \textbf{Regresión:} valor continuo (precio). \\
            \textbf{Clasificación:} categoría (spam/no spam).
        \end{exampleblock}
        \vspace{0.2cm}
        \begin{exampleblock}{No supervisado}
            Sin etiquetas. \\
            \textbf{Clustering:} agrupar clientes. \\
            \textbf{Reducción dimensionalidad:} PCA.
        \end{exampleblock}

        \column{0.48\textwidth}
        \begin{exampleblock}{Por refuerzo}
            Agente interactúa con entorno. \\
            Recompensas, maximización acumulada. \\
            Ej: robot aprendiendo a caminar.
        \end{exampleblock}
        \vspace{0.2cm}
        \begin{exampleblock}{Semi-supervisado / Autosupervisado}
            Combinan datos etiquetados y no etiquetados. \\
            Ej: predecir la siguiente palabra (NLP).
        \end{exampleblock}
    \end{columns}
    % IMAGEN: Esquema de los cuatro tipos
\end{frame}

% ============================================================================
% SECCIÓN 2: FUNDAMENTOS MATEMÁTICOS
% ============================================================================
\section{Fundamentos Matemáticos}

\begin{frame}{Estructuras de datos en IA}
    \begin{block}{Escalares, Vectores, Matrices y Tensores}
        \begin{itemize}
            \item \textbf{Escalar:} número real (bias, temperatura). Ej: $x = 5.2$
            \item \textbf{Vector:} lista ordenada (punto en espacio de características). Ej: $\mathbf{x}=[2.5,1.8,3.0]$
            \item \textbf{Matriz:} tabla 2D (imagen en escala de grises). Ej: $\mathbf{X}\in\mathbb{R}^{n\times m}$
            \item \textbf{Tensor:} generalización a más dimensiones (lote de imágenes a color). Ej: $\mathcal{T}\in\mathbb{R}^{B\times H\times W\times C}$
        \end{itemize}
    \end{block}
    % IMAGEN: Representación gráfica de las cuatro estructuras
\end{frame}

\begin{frame}{Funciones en IA}
    \begin{columns}[t]
        \column{0.5\textwidth}
        \begin{itemize}
            \item \textbf{Lineales:} $f(x)=mx+b$ (regresión)
            \item \textbf{Polinomiales:} $f(x)=a_nx^n+\dots$ (aproximación)
            \item \textbf{Exponenciales:} $e^x$ (softmax, pérdida)
        \end{itemize}
        \column{0.5\textwidth}
        \begin{itemize}
            \item \textbf{Logarítmicas:} $\ln x$ (entropía)
            \item \textbf{Trigonométricas:} $\sin$, $\cos$ (positional encodings)
        \end{itemize}
    \end{columns}
    % IMAGEN: Gráficas de estos tipos
\end{frame}

% ============================================================================
% SECCIÓN 3: FUNCIONES DE ACTIVACIÓN Y PÉRDIDA
% ============================================================================
\section{Funciones de Activación y Pérdida}

\begin{frame}{Funciones de activación (I)}
    Introducen no linealidad.
    \begin{columns}[t]
        \column{0.33\textwidth}
        \begin{block}{Sigmoide}
            $\sigma(x)=\frac{1}{1+e^{-x}}\in(0,1)$
        \end{block}
        \begin{block}{Tanh}
            $\tanh(x)\in(-1,1)$
        \end{block}
        \column{0.33\textwidth}
        \begin{block}{ReLU}
            $\max(0,x)$
        \end{block}
        \begin{block}{Leaky ReLU}
            $\max(\alpha x,x)$
        \end{block}
        \column{0.33\textwidth}
        \begin{block}{ELU}
            $\begin{cases}x&x\ge0\\ \alpha(e^x-1)&x<0\end{cases}$
        \end{block}
    \end{columns}
    % IMAGEN: Gráficas superpuestas
\end{frame}

\begin{frame}{Funciones de activación (II)}
    \begin{columns}[t]
        \column{0.5\textwidth}
        \begin{block}{Softmax}
            $\displaystyle \text{softmax}(\mathbf{z})_i = \frac{e^{z_i}}{\sum_{j=1}^K e^{z_j}}$ \\
            Convierte logits en distribución de probabilidad.
        \end{block}
        \column{0.5\textwidth}
        \begin{block}{GELU / Swish}
            GELU: $x\,\Phi(x)$ (usada en transformers). \\
            Swish: $x\cdot\sigma(x)$ (EfficientNet).
        \end{block}
    \end{columns}
    % IMAGEN: Gráficas de GELU y Swish
\end{frame}

\begin{frame}{Funciones de pérdida}
    \begin{columns}[t]
        \column{0.5\textwidth}
        \begin{alertblock}{Regresión}
            \begin{itemize}
                \item MSE: $\frac{1}{n}\sum (y_i-\hat{y}_i)^2$
                \item MAE: $\frac{1}{n}\sum |y_i-\hat{y}_i|$
            \end{itemize}
        \end{alertblock}
        \column{0.5\textwidth}
        \begin{alertblock}{Clasificación}
            \begin{itemize}
                \item Log-loss: $-\sum [y_i\log\hat{y}_i+(1-y_i)\log(1-\hat{y}_i)]$
                \item Hinge loss: $\max(0,1-y_i\hat{y}_i)$
                \item Focal loss (clases desbalanceadas)
            \end{itemize}
        \end{alertblock}
    \end{columns}
    % IMAGEN: Comparación de curvas de pérdida
\end{frame}

% ============================================================================
% SECCIÓN 4: LÓGICA Y NOTACIÓN
% ============================================================================
\section{Lógica y Notación}

\begin{frame}{Lógica proposicional en IA}
    \begin{block}{Aplicaciones}
        \begin{itemize}
            \item \textbf{Árboles de decisión:} nodos con condiciones lógicas.
            \item \textbf{Sistemas basados en reglas:} motores de inferencia.
            \item \textbf{Programación lógica} (Prolog).
            \item \textbf{Verificación de modelos} en sistemas críticos.
        \end{itemize}
    \end{block}
    % IMAGEN: Ejemplo de árbol de decisión
\end{frame}

\begin{frame}{Notación de sumatorias y productorias}
    \begin{columns}[t]
        \column{0.5\textwidth}
        \begin{exampleblock}{Sumatoria}
            $\displaystyle \sum_{i=1}^{n} x_i = x_1 + x_2 + \dots + x_n$ \\
            Ej: $z = \sum_{i=1}^{n} w_i x_i + b$
        \end{exampleblock}
        \column{0.5\textwidth}
        \begin{exampleblock}{Productoria}
            $\displaystyle \prod_{i=1}^{n} x_i = x_1 \cdot x_2 \cdots x_n$ \\
            Ej: $P(X_1,\dots,X_n) = \prod P(X_i)$
        \end{exampleblock}
    \end{columns}
\end{frame}

% ============================================================================
% SECCIÓN 5: HERRAMIENTAS DE PYTHON
% ============================================================================
\section{Herramientas de Python}

\begin{frame}{Librerías esenciales}
    \begin{columns}[t]
        \column{0.5\textwidth}
        \begin{block}{Fundamentales}
            NumPy, SciPy, Matplotlib, SymPy
        \end{block}
        \begin{block}{Machine Learning clásico}
            Scikit-learn, XGBoost, LightGBM, CatBoost
        \end{block}
        \column{0.5\textwidth}
        \begin{block}{Deep Learning}
            TensorFlow/Keras, PyTorch, JAX
        \end{block}
        \begin{block}{Especializadas}
            NLP: HuggingFace, NLTK, spaCy \\
            Visión: OpenCV, torchvision \\
            Audio: librosa \\
            Grafos: NetworkX, PyG
        \end{block}
    \end{columns}
    % IMAGEN: Logos de las principales librerías
\end{frame}

\begin{frame}[fragile]{Ejemplo práctico en Python}
    \begin{lstlisting}
import numpy as np
import matplotlib.pyplot as plt

def sigmoid(x): return 1/(1+np.exp(-x))
x = np.linspace(-5,5,100)
plt.plot(x, sigmoid(x))
plt.title('Función Sigmoide')
plt.grid(True)
plt.show()
    \end{lstlisting}
    % IMAGEN: Gráfica generada
\end{frame}

% ============================================================================
% SECCIÓN 6: MODELOS DE MACHINE LEARNING
% ============================================================================
\section{Modelos de Machine Learning}

\subsection{Supervisados}

\begin{frame}{Regresión Lineal}
    \begin{block}{Formulación}
        $y = \beta_0 + \beta_1 x_1 + \dots + \beta_p x_p + \varepsilon$
    \end{block}
    \begin{exampleblock}{Función de coste (MSE)}
        $J(\beta) = \frac{1}{n}\sum (y_i - \hat{y}_i)^2$
    \end{exampleblock}
    \begin{itemize}
        \item Solución analítica: $\hat{\beta} = (X^TX)^{-1}X^Ty$ (coste $O(p^3)$)
        \item Gradiente descendente: $\nabla J = \frac{2}{n}X^T(X\beta - y)$
        \item \textbf{Supuestos:} linealidad, homocedasticidad, normalidad (para inferencia)
    \end{itemize}
    % Basado en Pregunta 3 del solucionario
    % IMAGEN: Recta de regresión con puntos
\end{frame}

\begin{frame}{Regresión Logística}
    \begin{block}{Modelo de clasificación binaria}
        $P(y=1|x) = \sigma(\beta^T x) = \frac{1}{1+e^{-\beta^T x}}$
    \end{block}
    \begin{alertblock}{Log-loss (entropía cruzada)}
        $J(\beta) = -\sum [y_i\log\hat{y}_i + (1-y_i)\log(1-\hat{y}_i)]$
    \end{alertblock}
    \begin{itemize}
        \item Gradiente: $\nabla J = \sum (\hat{y}_i - y_i)x_i$
        \item No tiene solución cerrada; se optimiza con gradiente descendente.
    \end{itemize}
    % Basado en Pregunta 4
    % IMAGEN: Curva sigmoide con puntos clasificados
\end{frame}

\begin{frame}{k-Vecinos más cercanos (k-NN)}
    \begin{block}{Funcionamiento}
        Clasifica un nuevo punto por voto mayoritario de sus $k$ vecinos más cercanos (según distancia euclidiana, Manhattan, etc.)
    \end{block}
    \begin{itemize}
        \item Entrenamiento: $O(1)$ (solo almacena)
        \item Inferencia: $O(n\cdot p)$ (costoso con millones de datos)
        \item Elección de $k$: validación cruzada.
        \item $k$ pequeño $\rightarrow$ sobreajuste; $k$ grande $\rightarrow$ subajuste.
    \end{itemize}
    % Basado en Pregunta 6
    % IMAGEN: Clasificación con k-NN en 2D
\end{frame}

\begin{frame}{Naive Bayes}
    \begin{block}{Teorema de Bayes}
        $P(C_k|x) \propto P(C_k) \prod_{j=1}^n P(x_j|C_k)$
    \end{block}
    \begin{itemize}
        \item Supuesto de independencia condicional (``naive'').
        \item Para características continuas: $P(x_j|C_k) \sim \mathcal{N}(\mu_{kj},\sigma_{kj}^2)$
        \item Muy eficiente en problemas de texto (clasificación de documentos).
    \end{itemize}
    % Basado en Pregunta 7
    % IMAGEN: Esquema de probabilidades condicionales
\end{frame}

\begin{frame}{Máquinas de Vectores de Soporte (SVM)}
    \begin{block}{Idea principal}
        Maximizar el margen entre clases mediante un hiperplano óptimo.
    \end{block}
    \begin{itemize}
        \item \textbf{Vectores de soporte:} puntos que definen el margen.
        \item \textbf{Truco del kernel:} transformación implícita a mayor dimensión (lineal, RBF, polinomial).
        \item Parámetro $C$: controla penalización por errores.
    \end{itemize}
    % Basado en Preguntas 8 y 9
    % IMAGEN: SVM con margen y vectores
\end{frame}

\begin{frame}{Árboles de Decisión}
    \begin{block}{Partición recursiva}
        Cada nodo interno pregunta sobre una característica, maximizando la pureza de los hijos.
    \end{block}
    \begin{columns}[t]
        \column{0.5\textwidth}
        \textbf{Índice de Gini:} $1-\sum p_k^2$
        \column{0.5\textwidth}
        \textbf{Entropía:} $-\sum p_k\log p_k$
    \end{columns}
    \vspace{0.2cm}
    \begin{itemize}
        \item Interpretables, pero tienden a sobreajustar.
        \item Mitigación: poda, profundidad máxima, mínimo de muestras por hoja.
    \end{itemize}
    % Basado en Pregunta 10
    % IMAGEN: Árbol de decisión
\end{frame}

\begin{frame}{Random Forest}
    \begin{block}{Ensemble de árboles}
        Construye múltiples árboles con muestras bootstrap (bagging) y aleatoriedad en características.
    \end{block}
    \begin{itemize}
        \item Reduce la varianza sin aumentar el sesgo.
        \item Predicción por promedio (regresión) o voto mayoritario (clasificación).
        \item Robusto y de buen rendimiento en problemas tabulares.
    \end{itemize}
    % Basado en Pregunta 11
    % IMAGEN: Esquema de bagging
\end{frame}

\begin{frame}{Redes Neuronales (MLP)}
    \begin{block}{Arquitectura}
        Capas de neuronas: entrada, ocultas, salida. Cada neurona: $a = \sigma(\sum w_j x_j + b)$.
    \end{block}
    \begin{itemize}
        \item \textbf{Backpropagation:} gradientes por regla de la cadena.
        \item Activaciones no lineales esenciales para aproximar funciones complejas.
    \end{itemize}
    % Basado en Pregunta 12
    % IMAGEN: Diagrama de una neurona
\end{frame}

\subsection{No Supervisados}

\begin{frame}{K-Means}
    \begin{block}{Objetivo}
        Particionar datos en $k$ clusters minimizando la suma de distancias cuadráticas intra-cluster:
        $J = \sum_{j=1}^k \sum_{x\in C_j} \|x-\mu_j\|^2$
    \end{block}
    \begin{itemize}
        \item Algoritmo iterativo: asignación y actualización de centroides.
        \item Elección de $k$: método del codo, silhouette score.
    \end{itemize}
    % Basado en Pregunta 13
    % IMAGEN: Clusters K-Means
\end{frame}

\begin{frame}{PCA (Análisis de Componentes Principales)}
    \begin{block}{Reducción de dimensionalidad}
        Proyecta datos a un nuevo espacio de menor dimensión conservando la máxima varianza.
    \end{block}
    \begin{itemize}
        \item Componentes principales = autovectores de la matriz de covarianza.
        \item Se conservan los primeros $k$ componentes que explican un porcentaje de varianza (ej. 95\%).
    \end{itemize}
    % Basado en Pregunta 14
    % IMAGEN: Proyección 2D a 1D
\end{frame}

\begin{frame}{DBSCAN}
    \begin{block}{Clustering basado en densidad}
        Agrupa puntos que están estrechamente empaquetados, marcando como ruido los puntos aislados.
    \end{block}
    \begin{itemize}
        \item \textbf{Puntos núcleo, borde y ruido.}
        \item Detecta clusters de forma arbitraria y outliers.
        \item Parámetros: \texttt{eps} (radio) y \texttt{minPts}.
    \end{itemize}
    % Basado en Pregunta 15
    % IMAGEN: Clusters no esféricos con ruido
\end{frame}

\subsection{Aprendizaje por Refuerzo}

\begin{frame}{Conceptos básicos de RL}
    \begin{block}{Elementos}
        \textbf{Agente, entorno, estado, acción, recompensa.}
    \end{block}
    \begin{itemize}
        \item Objetivo: maximizar recompensa acumulada $G_t = \sum_{k=0}^\infty \gamma^k R_{t+k+1}$.
        \item \textbf{Ecuación de Bellman:} $Q(s,a) = r + \gamma \max_{a'} Q(s',a')$.
        \item \textbf{Q-learning:} $Q(s,a) \leftarrow Q(s,a) + \alpha [r + \gamma \max_{a'} Q(s',a') - Q(s,a)]$.
        \item \textbf{Exploración vs explotación} (ej. $\epsilon$-greedy).
    \end{itemize}
    % Basado en Pregunta 16
    % IMAGEN: Diagrama agente-entorno
\end{frame}

% ============================================================================
% SECCIÓN 7: EVALUACIÓN Y METODOLOGÍA
% ============================================================================
\section{Evaluación y Metodología}

\begin{frame}{Métricas de evaluación}
    \begin{columns}[t]
        \column{0.5\textwidth}
        \begin{block}{Clasificación}
            \begin{itemize}
                \item Accuracy, precisión, recall, F1, AUC-ROC
                \item Precisión = $\frac{VP}{VP+FP}$
                \item Recall = $\frac{VP}{VP+FN}$
            \end{itemize}
        \end{block}
        \column{0.5\textwidth}
        \begin{block}{Regresión}
            \begin{itemize}
                \item MSE, MAE, RMSE, $R^2$
            \end{itemize}
        \end{block}
    \end{columns}
    \vspace{0.2cm}
    \begin{exampleblock}{Validación cruzada}
        Divide los datos en $k$ pliegues; más robusta que una única partición.
    \end{exampleblock}
    % Basado en Pregunta 18
    % IMAGEN: Matriz de confusión
\end{frame}

\begin{frame}{Metodología CRISP-DM}
    \begin{block}{Fases}
        \begin{enumerate}
            \item Comprensión del negocio
            \item Comprensión de los datos
            \item Preparación de los datos
            \item Modelado
            \item Evaluación
            \item Despliegue
        \end{enumerate}
    \end{block}
    \begin{itemize}
        \item Proceso iterativo, no lineal.
        \item Aplicado a predicción de precios de viviendas (Pregunta 19).
    \end{itemize}
    % IMAGEN: Diagrama cíclico de CRISP-DM
\end{frame}

% ============================================================================
% SECCIÓN 8: CASO INTEGRADOR
% ============================================================================
\section{Caso integrador: Sistema de Recomendación}

\begin{frame}{Diseño de un sistema de recomendación (Netflix)}
    \begin{block}{Combinación de aprendizajes (Pregunta 20)}
        \begin{itemize}
            \item \textbf{Supervisado:} predecir rating o probabilidad de ver.
            \item \textbf{No supervisado:} segmentar usuarios (clustering).
            \item \textbf{Refuerzo contextual:} optimizar recomendaciones secuenciales (bandits).
        \end{itemize}
    \end{block}
    \begin{columns}[t]
        \column{0.5\textwidth}
        \textbf{Métricas offline:} Precision@k, Recall@k, NDCG.
        \column{0.5\textwidth}
        \textbf{Métricas online:} A/B test (tiempo de visualización, retención).
    \end{columns}
    \vspace{0.2cm}
    \textbf{Cold start:} para nuevos usuarios o contenido: popularidad, metadatos.
    % IMAGEN: Esquema de recomendación
\end{frame}

% ============================================================================
% SECCIÓN 9: CONCLUSIÓN
% ============================================================================
\section{Conclusión}

\begin{frame}{Resumen de la sesión}
    \begin{exampleblock}{Lo que hemos visto}
        \begin{itemize}
            \item Jerarquía IA $\supset$ ML $\supset$ DL y tipos de aprendizaje.
            \item Fundamentos matemáticos: vectores, matrices, funciones.
            \item Funciones de activación y pérdida.
            \item Herramientas de Python (NumPy, PyTorch, etc.).
            \item Modelos supervisados, no supervisados y por refuerzo.
            \item Evaluación y metodología CRISP-DM.
            \item Caso integrador de recomendación.
        \end{itemize}
    \end{exampleblock}
    \vspace{0.3cm}

\end{frame}

% --- Diapositiva final ---
\begin{frame}
    \centering
    \Huge \textbf{¡Gracias!}
    \vspace{1cm}

\end{frame}

\end{document}